\section{Round-twenty-nine reconstruction: Jacobian-weighted grammars and mixed local limits}
\label{sec:r29-a2}

The previous certificate contained an unweighted sum of inverse-branch
$C^2$ norms and was therefore empty at derivative order zero.  The present
certificate weights every branch derivative by the inverse Jacobian, derives
the bad-word estimate from a finite automaton, and uses enough continuous roof
generators to have trivial annihilator in dimensions one and two.

\subsection{A finite geometric grammar}

Let $T_a:M_a\to M_a$ be a compact $C^{r_*+5}$ family of finite-horizon
dispersing billiards with strictly positive curvature and no corners.  Fix a
finite Markov magnet $Y_a$, transported to one reference magnet $Y$.  Its
first-return inverse branches are denoted by $h$, with return time $R(h)$,
inverse unstable Jacobian $J_h$, homology/return displacement
$\kappa(h)\in\Z^3$, and continuous roof vector $\tau_h\in\R^{d_c}$,
$1\le d_c\le2$.

A \emph{Jacobian-weighted finite grammar certificate} consists of:

\begin{enumerate}[label=(G\arabic*)]
\item a finite family of parent-fold chart types and constants such that for
      $0\le j\le r_*+3$,
\[
 \sum_{h:R(h)=n}\sup_{|\alpha|\le j+2}
 \norm{D^\alpha(J_h,\,J_h\tau_h,\,J_h\partial_a^jh)}_{L^1(Y)}
 \le C_j(1+n)^{q_j}e^{-cn};                               \tag{A2.1}
\]
\item a finite directed automaton whose paths encode all return words, with a
      distinguished collection of returned-UNI blocks; every path either
      contains two disjoint distinguished blocks or ends in one of finitely
      many terminal nonstationary chart types;
\item for each distinguished block, two matched inverse branches on the same
      stable-curve chart and a map
      $\Psi=(\Psi_1,\ldots,\Psi_{d_c})$ satisfying
\[
 \abs{\det D_z\Psi}\ge c_{\rm uni}>0                     \tag{A2.2}
\]
      on a boundary-flat subchart, with all derivatives through order $r_*$
      controlled by (A2.1);
\item for each terminal chart type, a coordinate $z_1$ with
\[
 \abs{\partial_{z_1}(b\cdot\tau_h)}
 \ge c_{\rm ns}|b|(1+R(h))^{-q}                           \tag{A2.3}
\]
      whenever $|b|\ge1$;
\item periodic orbit differences whose lattice coordinates generate $\Z^3$
      and continuous roof differences
      $\beta_0,\ldots,\beta_{d_c}\in\R^{d_c}$ satisfying
\[
 \{b\in\R^{d_c}:b\cdot\beta_j\in2\pi\Z
                 \text{ for all }j\}=\{0\}.              \tag{A2.4}
\]
\end{enumerate}

No item above is an operator spectral estimate.  Each item is a finite chart,
finite automaton, weighted derivative sum, or finite periodic datum.

\begin{proposition}[Certificate validation and nonemptiness]
\label{prop:r29-a2-certificate}
The certificate is stable under sufficiently small $C^{r_*+5}$ perturbations.
Its class is nonempty.  In dimension $d_c=1$ the continuous generators may be
chosen as $1$ and $\sqrt2$; in dimension $d_c=2$ they may be chosen as
\[
 \beta_0=(1,0),\qquad \beta_1=(0,1),\qquad
 \beta_2=(\sqrt2,\sqrt3).                                \tag{A2.5}
\]
A finite-horizon circular Lorentz table can be perturbed in three pairwise
disjoint regular orbit tubes so that its corresponding roof differences have
these values while preserving the finite Markov grammar and (A2.1)--(A2.3).
\end{proposition}

\begin{proof}
For a finite-horizon dispersing table, the growth lemma and bounded distortion
on homogeneity strips give exponential return tails.  Differentiating the
inverse-branch formula always leaves the factor $J_h$; summing over a return
layer therefore gives (A2.1), including at $j=0$.  A finite subcover of the
magnet and its first-return singularity partition gives the chart-type
automaton.  Transversality of two returned flight-time differences is open;
if a terminal word misses both transverse blocks, the final regular flight
has a nonstationary impact coordinate, giving (A2.3) after a finite
refinement.

To realize the roof data, choose three disjoint regular periodic tubes.  A
normal displacement of one scatterer arc changes only the flight length in
its tube to first order, with nonzero derivative.  The finite-dimensional
implicit-function theorem therefore prescribes the three differences in
(A2.5).  Blocking arcs may be added outside the tubes and then smoothed,
preserving positive curvature and finite horizon.  All inequalities are
strict and hence survive a small perturbation.

For $d_c=2$, if $b\cdot\beta_j\in2\pi\Z$ for all three generators, the first
two identities give $b=2\pi(m,n)$.  The third gives
$m\sqrt2+n\sqrt3\in\Z$, which forces $m=n=0$ by the rational independence of
$1,\sqrt2,\sqrt3$.  The one-dimensional argument is identical with
$1,\sqrt2$.  Thus (A2.4) holds.
\end{proof}

\subsection{Completed anisotropic spaces}

Let $\mathcal W^s$ be the compact family of homogeneous stable curves of
length at most $\delta_0$.  For $0<\alpha<1$, integer $r_*$, and a smooth
function $f$, set
\begin{equation}\tag{A2.6}
\begin{split}
 \norm{f}_{\cB^{r_*}}={}&
 \sum_{j=0}^{r_*}\sup_{W\in\mathcal W^s}
 \sup_{\norm{\varphi}_{C^{r_*-j+\alpha}(W)}\le1}
 \abs{\int_WD_u^jf\,\varphi\,dm_W}\\
 &+\sup_{W,W'}{\abs{\int_Wf\varphi-\int_{W'}f\varphi' }
             \over d_{\mathcal W}(W,W')^\alpha},          
\end{split}
\end{equation}
where matched tests are transported by stable holonomy and homogeneity
weights are included in $d_{\mathcal W}$.  Let $\cB^{r_*}$ be the completion
and let $\cB_w$ be the same construction with $j=0$ and exponent $\alpha/2$.
The embedding $\cB^{r_*}\hookrightarrow\cB_w$ is compact because stable
curves and tests are compact after the homogeneity weights are imposed.

For $\xi=(u,b)\in[-\pi,\pi]^3\times\R^{d_c}$ define
\[
 \mathcal L_\xi f
 =\sum_hJ_h e^{iu\cdot\kappa(h)+ib\cdot\tau_h}f\circ h.   \tag{A2.7}
\]
Equation (A2.1) proves boundedness and $C^{r_*}$ source/parameter
differentiability from $\cB^{r_*}$ to $\cB_w$.

\subsection{Returned cancellation and finite bad-word closure}

\begin{lemma}[Matched returned-UNI block]
\label{lem:r29-a2-uni}
For a word containing two disjoint returned-UNI blocks, every stable-curve
pairing of its branch contribution is a finite sum of integrals
\[
 \int_Ue^{ib\cdot\Psi(z)}A_{W,\varphi}(z)\,dz,             \tag{A2.8}
\]
where, for $M\le r_*$,
\[
 \sum\sum_{|\alpha|\le M}
\norm{D^\alpha A_{W,\varphi}}_{L^1(U)}
 \le C_Me^{-c n}\norm{f}_{\cB^{r_*}}
\norm{\varphi}_{C^M}. \tag{A2.9}
\]
Hence its strong-to-weak contribution is bounded by
$C_Me^{-cn}(1+|b|)^{-M}$, uniformly after two source derivatives.
\end{lemma}

\begin{proof}
Matched branches are first pulled through their common suffix, so the
oscillatory phase is the returned temporal difference and both amplitudes live
on the same original stable chart.  The boundary-flat partition eliminates
boundary distributions.  The vector fields
$V_j=\sum_k(D\Psi)^{-1}_{jk}\partial_{z_k}$ satisfy
$V_j\Psi_k=\delta_{jk}$.  Repeated integration by parts transfers $M$
derivatives to the amplitude.  Every derivative is one of the weighted terms
in (A2.1), and the matched-curve difference is controlled by the last term of
(A2.6).  Summing chart types gives (A2.9).
\end{proof}

\begin{lemma}[Finite-automaton bad-word estimate]
\label{lem:r29-a2-bad}
Let $\mathfrak B_n$ be the words with fewer than two distinguished blocks.
Then for every $M\le r_*$,
\[
 \norm{\sum_{\mathcal W\in\mathfrak B_n}
       \mathcal L_{\xi,\mathcal W}}_{\cB^{r_*}\to\cB_w}
 \le C_M(1+n)^{q_M}e^{-cn}(1+|b|)^{-M}.                  \tag{A2.10}
\]
\end{lemma}

\begin{proof}
The bad words form a regular language recognized by the finite certificate
automaton.  Its weighted transition matrix has spectral radius strictly below
that of the full return matrix because every recurrent component reaches a
distinguished block.  This gives $C(1+n)^qe^{-cn}$ total Jacobian weight.
Every terminal state carries one of the finitely many coordinates in (A2.3).
Integration by parts there gives $(1+|b|)^{-M}$; the polynomial inverse
nonstationarity factor is absorbed by the exponential return tail in (A2.1).
This derives the bad-word estimate from finite data instead of postulating an
infinite condition on every word.
\end{proof}

\subsection{Arithmetic, intermediate frequencies, and the local theorem}

Let $\mathscr G$ be the closed subgroup generated by the periodic differences
of $(\kappa,\tau,n)$.

\begin{lemma}[Correct continuous-character annihilator]
\label{lem:r29-a2-annihilator}
The periodic data in (G5) imply
\[
 \mathscr G=\Z^3\times\R^{d_c}\times\Z.                  \tag{A2.11}
\]
Therefore $\mathcal L_\xi$ has no unit-modulus eigenvalue for
$\xi\ne0$ modulo the lattice torus, and the centred covariance is positive
definite on the quotient by exact conservation relations.
\end{lemma}

\begin{proof}
The lattice periodic differences kill all lattice characters.  A continuous
character $x\mapsto e^{ib\cdot x}$ annihilates the continuous roof subgroup
exactly when the condition on the left of (A2.4) holds; hence $b=0$.  Coprime
periods kill the remaining constant phase.  Pontryagin duality gives
(A2.11).  A unit-modulus eigenfunction yields the same periodic character
relation.  A zero covariance direction is a square-integrable coboundary and
therefore has zero periodic sums, again contradicting the certificate.
\end{proof}

\begin{theorem}[All-frequency anisotropic estimate]
\label{thm:r29-a2-fourier}
Choose $M>d_c+8$ and $r_*\ge M+5$.  There are
$\delta,B,c,C,\varkappa>0$ such that
\[
 \norm{\mathcal L_\xi^n}_{\cB^{r_*}\to\cB_w}
 \le
 \begin{cases}
 C e^{-cn|\xi|^2}+Ce^{-cn},&|\xi|\le\delta,\\
 Ce^{-cn},&\delta\le|\xi|\le B,\\
 C\exp[-cn/(1+\log|b|)^q],&B<|b|\le e^{\varkappa n},\\
 C(1+n)^C(1+|b|)^{-M},&|b|>e^{\varkappa n}.
 \end{cases}                                             \tag{A2.12}
\]
The same bound holds after two source derivatives and is integrable in the
continuous Fourier variable.
\end{theorem}

\begin{proof}
Near zero, analytic perturbation of the simple eigenvalue gives the Gaussian
quadratic form and a contracting complement.  On the compact annulus,
Lemma~\ref{lem:r29-a2-annihilator}, compact embedding, and a contradiction
argument give a uniform gap.  For intermediate frequencies, divide the word
into blocks of length $C(1+\log|b|)^q$.  The grammar forces a positive fraction
of blocks either to contain a returned-UNI pair or to enter the exponentially
small bad automaton.  On a good block, Lemma~\ref{lem:r29-a2-uni} gives a fixed
cone loss in the actual stable-curve pairing; multiplying the losses gives the
third line.  At high frequency use Lemmas~\ref{lem:r29-a2-uni} and
\ref{lem:r29-a2-bad} directly.  The chosen derivative budget makes the last
line integrable and stable under source differentiation.
\end{proof}

Let $Z_n=(S_n\kappa,S_n\tau)$, with counting measure in the lattice
coordinates and Lebesgue measure in the continuous coordinates.

\begin{theorem}[Source-uniform mixed local limit theorem]
\label{thm:r29-a2-llt}
Uniformly for central deviations $(k_n,y_n)=O(\sqrt n)$ and sources in a fixed
$C^2$ ball,
\[
 P_H\{S_n\kappa=k_n,\ S_n\tau\in dy_n\}
 ={\exp[-\frac1{2n}(z_n-nm_H)^T\Sigma_H^{-1}(z_n-nm_H)]
   \over(2\pi n)^{(3+d_c)/2}\sqrt{\det\Sigma_H}}
 \left(1+O(n^{-1/2})\right)dy_n,                          \tag{A2.13}
\]
where $z_n=(k_n,y_n)$ and the lattice covolume is included in the numerator.
The error and its first two source derivatives are uniform.  Aperiodic lattice
classes outside the generated subgroup have probability zero.
\end{theorem}

\begin{proof}
Fourier inversion is split according to (A2.12).  The central ball uses the
third-order eigenvalue expansion and gives (A2.13).  The compact and
intermediate regions are exponentially or superpolynomially small.  The
high-frequency region is absolutely integrable because $M>d_c$.  The same
majorant controls two source derivatives.  The annihilator theorem identifies
the exact lattice support and removes spurious continuous characters.
\end{proof}

The certificate is now nonempty, Jacobian weighted, finite at the bad-word
level, arithmetically correct for $d_c=2$, and sufficient for a genuinely
operator-valued all-frequency proof.
